%============================================================================
% LICENSING NOTICE
%============================================================================
% This WIP file, upon integration into guide-v.tex, becomes part of the
% TMS/DMS/CMS Usage Guide and is released under CC BY-NC 4.0.
% For details, see LICENSE in the repository root.
%============================================================================

%============================================================================
% FALCON BMS TMS/DMS/CMS HOTAS GUIDE
% WIP FILE TEMPLATE V1.0 — FINAL (Briefing)
%============================================================================

% IMPORTANTE: Este é um TEMPLATE padronizado para arquivos WIP (Work-In-Progress)
% que serão integrados ao guide.tex.

% Nomenclatura: Siga wip-naming-v1.4
% Padrão: section-C{N}-S{M}[-S{K}]-{titulo}-{status}-{data}.tex
% Exemplo: section-C5-S2-cms-actuation-hotas-tables-final-2026-01-10.tex

% Status: dev | review | final | approved | deprecated
% Locação: WIP/ (ativo) | ARCHIVE/ (aprovado/descartado)

%============================================================================
% PREAMBLE COMPLETO — TMS/DMS/CMS Usage Guide for Falcon BMS 4.38.1
% Gerado: 11 February 2026
% Status: Improvements as per GitHub Issues #34 and #33
%============================================================================

\documentclass[11pt, a4paper, twoside]{report}

% --------------------------------------------------------------------------
% BASIC ENCODING AND LANGUAGE
% --------------------------------------------------------------------------

\usepackage[utf8]{inputenc}
\usepackage[T1]{fontenc}
\usepackage[english]{babel}

% --------------------------------------------------------------------------
% FONTS AND MICROTYPOGRAPHY
% --------------------------------------------------------------------------

\usepackage{lmodern}
\usepackage{microtype}

% --------------------------------------------------------------------------
% PAGE GEOMETRY AND LAYOUT
% --------------------------------------------------------------------------

\usepackage{geometry}
\geometry{a4paper, left=2.0cm, right=2.0cm, top=2.5cm, bottom=2.5cm}
\usepackage{setspace}
\onehalfspacing

% --------------------------------------------------------------------------
% COLORS AND LINKS AND PDF SETTINGS
% --------------------------------------------------------------------------

\usepackage[table]{xcolor}
\definecolor{linkblue}{HTML}{004488}
\definecolor{linkred}{HTML}{882222}
\definecolor{headerblue}{HTML}{003366}
\definecolor{rowgray}{HTML}{F5F5F5}
\definecolor{subheadgray}{HTML}{E0E0E0}
\usepackage{soul}
\usepackage[pdfencoding=auto, psdextra, colorlinks=true, linkcolor=linkblue, citecolor=linkred, urlcolor=linkblue, breaklinks=true]{hyperref}
\hypersetup{
	pdftitle={TMS, DMS and CMS Usage Guide for Falcon BMS},
	pdfauthor={Carlos "Metal" Nader},
	pdfsubject={Flight Simulation - Falcon BMS HOTAS Reference},
	pdfkeywords={Falcon BMS, F-16, HOTAS, TMS, DMS, CMS, Flight Simulation},
	pdfcreator={pdfLaTeX (MiKTeX)},
	pdfproducer={MiKTeX},
	pdflang={en-US},
}
\usepackage{bookmark}

% --------------------------------------------------------------------------
% CAPTION SETUP BY TYPE
% --------------------------------------------------------------------------

\usepackage{caption}

% GLOBAL PATTERN
\captionsetup{font=small, labelfont=bf, justification=centering, singlelinecheck=true}
% FIGURES
\captionsetup[figure]{font=footnotesize, labelfont=bf, justification=centering, singlelinecheck=true}
% % TABLE/LONGTABLE/hotastable:
\captionsetup[table]{font=small, labelfont=bf, justification=centering, singlelinecheck=true}

% --------------------------------------------------------------------------
% HEADERS AND FOOTERS (IMPROVED for report + twoside)
% --------------------------------------------------------------------------

\usepackage{fancyhdr}
\setlength{\headheight}{25pt}                    
\pagestyle{fancy}
\fancyhf{}                                        
\fancyhead[LO,RE]{\small\textit{\leftmark}}     
\fancyhead[RO,LE]{\small\thepage}               
\fancyfoot{}                                      
\renewcommand{\headrulewidth}{0.4pt}
\renewcommand{\footrulewidth}{0pt}

% --------------------------------------------------------------------------
% CHAPTER FORMATTING AND SPACING (via titlesec)
% --------------------------------------------------------------------------

\usepackage{titlesec}

% --------------------------------------------------------------------------
% CHAPTER - Nível 0 (Destaque Máximo)
% --------------------------------------------------------------------------
% Shape: display → standalone em nova linha
% Font: \Large\bfseries --- label e título ambos em bold para destaque máximo
% Numeração: SIM
% Espaçamento: 15pt antes, 28pt depois
% --------------------------------------------------------------------------

\titleformat{\chapter}[display]
{\normalfont\Large\bfseries}           % Formato: LARGE + bold (label)
{\chaptertitlename~\thechapter}        % Rótulo: "Chapter 1" (bold também)
{20pt}                                 % Espaço vertical antes do corpo
{\Large\bfseries}                      % Corpo do título: LARGE + bold

\titlespacing{\chapter}
{0pt}       % Margem esquerda (sem indent)
{15pt}      % Espaço ANTES do título do capítulo
{28pt}      % Espaço DEPOIS do título do capítulo (AUMENTADO)
[0pt]       % Margem direita

% --------------------------------------------------------------------------
% SECTION - Nível 1
% --------------------------------------------------------------------------
% Shape: hang → label à esquerda, corpo indentado (compacto)
% Font: \large\bfseries (mesmo destaque que chapter para coerência)
% Numeração: SIM
% Espaçamento: 15pt antes, 8pt depois (REDUZIDO de 15pt/10pt anteriores)
% --------------------------------------------------------------------------

\titleformat{\section}[hang]
{\normalfont\large\bfseries}           % Formato: large + bold
{\thesection}                          % Rótulo: "1", "2", etc
{1em}                                  % Espaço entre rótulo e corpo
{}                                     % Corpo do título (herda do format)

\titlespacing{\section}
{0pt}       % Margem esquerda (sem indent)
{15pt}      % Espaço ANTES do título da seção (REDUZIDO)
{8pt}       % Espaço DEPOIS do título da seção (REDUZIDO)
[0pt]       % Margem direita

% --------------------------------------------------------------------------
% SUBSECTION - Nível 2
% --------------------------------------------------------------------------
% Shape: hang → label esquerda, corpo indentado
% Font: \normalsize\bfseries (redução visual, ainda robusto)
% Numeração: SIM
% Espaçamento: 12pt antes, 6pt depois (compacto)
% --------------------------------------------------------------------------

\titleformat{\subsection}[hang]
{\normalfont\normalsize\bfseries}      % Formato: tamanho normal + bold
{\thesubsection}                       % Rótulo: "1.1", "1.2", etc
{1em}                                  % Espaço entre rótulo e corpo
{}                                     % Corpo do título

\titlespacing{\subsection}
{0pt}       % Margem esquerda (sem indent)
{12pt}      % Espaço ANTES do título da subseção
{6pt}       % Espaço DEPOIS do título da subseção
[0pt]       % Margem direita

% --------------------------------------------------------------------------
% SUBSUBSECTION - Nível 3
% --------------------------------------------------------------------------
% Shape: hang → label esquerda, máximo recuo (bem compacto)
% Font: \normalsize\bfseries (mesmo tamanho que subsection, padrão)
% Numeração: SIM (configurable com secnumdepth)
% Espaçamento: 10pt antes, 4pt depois (bem comprimido)
% --------------------------------------------------------------------------

\titleformat{\subsubsection}[hang]
{\normalfont\normalsize\bfseries}      % Formato: tamanho normal + bold
{\thesubsubsection}                    % Rótulo: "1.1.1", "1.1.2", etc
{1em}                                  % Espaço entre rótulo e corpo
{}                                     % Corpo do título

\titlespacing{\subsubsection}
{0pt}       % Margem esquerda (sem indent)
{10pt}       % Espaço ANTES do título
{4pt}       % Espaço DEPOIS do título (bem comprimido)
[0pt]       % Margem direita

% --------------------------------------------------------------------------
% PARAGRAPH - Nível 4 (CRÍTICO: Ilusão Óptica Resolvida)
% --------------------------------------------------------------------------
% Shape: runin → inline, integrado ao texto (seu requisito)
% Font: \small\bfseries (SOLUÇÃO PARA ILUSÃO ÓPTICA DO NEGRITO)%       
% Numeração: NÃO (padrão para \paragraph)
% Espaçamento: runin (sem espaço vertical antes, integrado ao texto)
% --------------------------------------------------------------------------

\titleformat{\paragraph}[runin]
{\normalfont\small\bfseries}           % Formato: SMALL + bold (compensa ilusão)
{}                                     % Rótulo: vazio (sem numeração)
{0em}                                  % Espaço entre rótulo e corpo (zero, runin)
{}                                     % Corpo do título

\titlespacing{\paragraph}
{0pt}       % Margem esquerda (sem indent)
{8pt}       % Espaço ANTES do título do parágrafo (ADICIONADO - cria separação)
{1em}       % Espaço DEPOIS do título (ADICIONADO - espaço após ":")
[0pt]       % Margem direita

% --------------------------------------------------------------------------
% SUBPARAGRAPH - Nível 5 (Preparação Futura)
% --------------------------------------------------------------------------
% Shape: runin → inline, integrado ao texto
% Font: \small (MENOS destaque que \paragraph, sem bold)
% Numeração: NÃO (padrão para \subparagraph)
% Espaçamento: runin (integrado)
% Uso: Se usado, aparecerá com destaque menor que \paragraph
% --------------------------------------------------------------------------

\titleformat{\subparagraph}[runin]
{\normalfont\small}                    % Formato: small, SEM bold (menos destaque)
{}                                     % Rótulo: vazio (sem numeração)
{0em}                                  % Espaço entre rótulo e corpo (zero, runin)
{}                                     % Corpo do título

\titlespacing{\subparagraph}
{0pt}       % Margem esquerda (sem indent)
{8pt}       % Espaço ANTES do título (ADICIONADO)
{1em}       % Espaço DEPOIS do título (ADICIONADO)
[0pt]       % Margem direita

% --------------------------------------------------------------------------
% TABLES AND MACROS
% --------------------------------------------------------------------------

\usepackage{booktabs}
\usepackage{array}
\usepackage{longtable}
\usepackage{tabularx}

% Custom Columns
\newcolumntype{L}[1]{>{\raggedright\arraybackslash}p{#1}}
\newcolumntype{C}[1]{>{\centering\arraybackslash}p{#1}}
\newcolumntype{R}[1]{>{\raggedleft\arraybackslash}p{#1}}

% Macro for Visual Reference Links
\newcommand{\imglink}[1]{\hspace{2pt}\hyperref[#1]{\scriptsize\textbf{[Fig]}}}

% --------------------------------------------------------------------------
% HOTAS LOT FILE HANDLE 
% --------------------------------------------------------------------------
\makeatletter
\newcommand{\listofhotastables}{%
	\section*{List of HOTAS Tables}%
	\addcontentsline{toc}{section}{List of HOTAS Tables}%
	\@starttoc{hotas}%
}
\makeatother

% --------------------------------------------------------------------------
% HOTAS table environment
% Version: 2.1 (2026-02-10)
% --------------------------------------------------------------------------

\newenvironment{hotastable}[1]{%
	\footnotesize
	\setlength{\tabcolsep}{3pt}
	\renewcommand{\arraystretch}{1.35}
	\begin{longtable}{L{1.00cm} L{0.90cm} L{0.90cm} L{3.30cm} L{6.40cm} L{1.40cm} L{1.60cm}}
		\caption{#1}\label{table.\thetable}\\
		\noalign{\addtocontents{hotas}{\protect\contentsline{table}{\protect\numberline{\thetable}#1}{\thepage}{table.\thetable}}}%
		\rowcolor{headerblue}
		\multicolumn{1}{>{\centering\arraybackslash}p{1.00cm}}{\textbf{\color{white}Mode}} &
		\multicolumn{1}{>{\centering\arraybackslash}p{0.90cm}}{\textbf{\color{white}Dir.}} &
		\multicolumn{1}{>{\centering\arraybackslash}p{0.90cm}}{\textbf{\color{white}Act.}} &
		\multicolumn{1}{>{\centering\arraybackslash}p{3.30cm}}{\textbf{\color{white}Function}} &
		\multicolumn{1}{>{\centering\arraybackslash}p{6.40cm}}{\textbf{\color{white}Effect / Nuance}} &
		\multicolumn{1}{>{\centering\arraybackslash}p{1.40cm}}{\textbf{\color{white}Dash34}} &
		\multicolumn{1}{>{\centering\arraybackslash}p{1.60cm}}{\textbf{\color{white}Train.}} \\
		\endfirsthead
		\rowcolor{headerblue}
		\multicolumn{1}{>{\centering\arraybackslash}p{1.00cm}}{\textbf{\color{white}Mode}} &
		\multicolumn{1}{>{\centering\arraybackslash}p{0.90cm}}{\textbf{\color{white}Dir.}} &
		\multicolumn{1}{>{\centering\arraybackslash}p{0.90cm}}{\textbf{\color{white}Act.}} &
		\multicolumn{1}{>{\centering\arraybackslash}p{3.30cm}}{\textbf{\color{white}Function}} &
		\multicolumn{1}{>{\centering\arraybackslash}p{6.40cm}}{\textbf{\color{white}Effect / Nuance}} &
		\multicolumn{1}{>{\centering\arraybackslash}p{1.40cm}}{\textbf{\color{white}Dash34}} &
		\multicolumn{1}{>{\centering\arraybackslash}p{1.60cm}}{\textbf{\color{white}Train.}} \\
		\endhead
		\multicolumn{7}{r}{\small\emph{Continued on next page}}\\
		\endfoot
		\endlastfoot
	}{%
	\end{longtable}
}

% --------------------------------------------------------------------------
% SIMPLE REFERENCE MACROS FOR BMS DOCS
% --------------------------------------------------------------------------

\newcommand{\dashref}[1]{Dash-34~\S~#1}
\newcommand{\dashone}[1]{Dash-1~\S~#1}
\newcommand{\trnref}[1]{TRN~#1}
\newcommand{\trnman}{BMS Training Manual 4.38.1}
\newcommand{\bmsver}{Falcon BMS~4.38.1}
\newcommand{\dashrefs}[1]{\textit{TO 1F-16CMAM-34-1-1}, Dash-34, sections \texttt{#1}}

% --------------------------------------------------------------------------
% CROSS-REFERENCE MACROS (INTERNAL GUIDE)
% --------------------------------------------------------------------------

\newcommand{\secref}[1]{\hyperref[#1]{Section~\ref*{#1}}}
\newcommand{\chapref}[1]{\hyperref[#1]{Chapter~\ref*{#1}}}
\newcommand{\tabref}[1]{\hyperref[#1]{Table~\ref*{#1}}}
\newcommand{\figref}[1]{\hyperref[#1]{Figure~\ref*{#1}}}

% --------------------------------------------------------------------------
% VERSION CONTROL MACROS
% --------------------------------------------------------------------------

\newcommand{\docversion}{future v0.5.0.0}
\newcommand{\docbuild}{16 February 2026}
\newcommand{\docstartdate}{05 January 2026}
\newcommand{\docenddate}{}
\newcommand{\chapterscompletedof}{4/7}
\newcommand{\tablesfilledpct}{}
\newcommand{\fulldocversion}{\docversion+\docbuild}

% --------------------------------------------------------------------------
% GRAPHICS
% --------------------------------------------------------------------------

\usepackage{graphicx}
\graphicspath{{fig/}}
\usepackage{float}
\usepackage{enumitem}

% --------------------------------------------------------------------------
% TITLE
% --------------------------------------------------------------------------

\title{TMS, DMS and CMS Usage Guide for \bmsver}
\author{Carlos ``Metal'' Nader}
\date{Version \fulldocversion{} | Progress: Chapters \chapterscompletedof{} | Tables \tablesfilledpct{} | January 2026}

% ============================================================================
% DOCUMENT BEGIN
% ============================================================================

\begin{document}
	
	\maketitle
	
	\pagenumbering{roman}
	
	% --------------------------------------------------------------------------
	% TOC DEPTH CONFIGURATION
	% --------------------------------------------------------------------------
	\setcounter{tocdepth}{3}       % Show up to \subsubsection in TOC
	\setcounter{secnumdepth}{3}    % Number up to \subsubsection
	
	\newpage
	\tableofcontents
	\newpage
	
	% --------------------------------------------------------------------------
	% LIST OF HOTAS TABLES
	% --------------------------------------------------------------------------
	\phantomsection
	\listofhotastables
	\newpage
	
	\pagenumbering{arabic}

%============================================================================
% WIP FILE METADATA (NOT RENDERED IN PDF)
%============================================================================
% File Name: section-C4-S3-tms-aa-review-2026-02-17.tex
% WIP Naming Convention: v1.4
% Target Chapter: C4 (TMS)
% Target Section: S3 (TMS in Air-to-Air)
% Target Subsection: S5 (TMS and IFF Interrogation)
%
% WIP Status: review
% Created: 2026-02-13
% Last Modified: 2026-02-17
% Integration Status: NOT YET INTEGRATED | INTEGRATION TARGET: v0.5.0.0
%
% Narrative Completion: ~28.57% (C4:S3 intro complete, C4:S3:S5 intro complete, C4:S3:S5:S1 complete, C4:S3:S5:S2 complete)
% Table Fill Status: 0% (0 of 6 tables populated)
%
% Completion Status by Subsection:
% - C4:S3 Introduction (main section): FINAL ✓
% - C4:S3:S1 (SAM): Draft structure only
% - C4:S3:S2 (STT): Draft structure only
% - C4:S3:S3 (TWS): Draft structure only
% - C4:S3:S4 (ACM): Draft structure only
% - C4:S3:S5 (IFF): Introduction FINAL ✓ | S1 FINAL ✓ | S2 FINAL ✓ | Subsubsections S3-S5: Draft structure only
% - C4:S3:S6 (TGP A-A): Draft structure only
%
% Table Status (0 of 6 complete):
% - Table C4-S3-SAM: Placeholder only
% - Table C4-S3-STT: Placeholder only
% - Table C4-S3-TWS: Placeholder only
% - Table C4-S3-ACM: Placeholder only
% - Table C4-S3-IFF: Placeholder only
% - Table C4-S3-TGP: Placeholder only
%
% Dependencies:
% - Requires: section-C2-S1-S3-soi-context-dependency-final-2026-02-15.tex
%   (TGP introduction in C2 fundamentals, integration target: v0.5.0.0)
% - Requires: C4:S2:S1 (HSD/PDLT) for cross-reference in Note
%   (currently placeholder, to be developed)
%
% Changes This Revision (2026-02-17):
% - C4:S3:S5 Introduction: clarified "cockpit identification methods" (vs ambiguous "identification procedures")
% - C4:S3:S5:S1 (SCAN Interrogation Mode): FINAL ✓
%   * Parágrafo 1: SCAN overview + TMS Left short press (<0.6s) bold + 32 reply limit
%   * \paragraph{Configuration:}: DED INTG page + FCR CNTL OSB 10 + decouple/couple logic
%   * Cross-refs validated: \dashref{2.2.4.4.2.1.1}, \dashref{2.3.1.5.1.6.2}
%   * Parágrafo 2: Symbology (green circle = friendly, yellow square = bogey/unknown)
%   * Parágrafo 3: Tactical context (BVR intercepts, ROE compliance, multiple contacts)
%   * Corrected tactical use paragraph (removed erroneous LOS content)
% - C4:S3:S5:S2 (LOS Interrogation Mode): FINAL ✓
%   * Parágrafo 1: Concept (beam coverage area, cursor-directed, 2 beams, vs SCAN comparison)
%   * Parágrafo 2: Behavior (TMS Left ≥0.6s, 16 replies, beam geometry nuance, range filtering)
%   * \paragraph{Configuration:}: DED INTG + DCS right (SEQ to LOS page), independent from SCAN
%   * Cross-refs validated: \dashref{2.2.4.4.2.1}, \dashref{2.2.4.4.2.1.2}, \dashref{2.2.4.4.2.2}
%   * Parágrafo 3: Symbology (cross-ref to S1, reduced redundancy)
%   * Parágrafo 4: Tactical use (removed unsupported "VID passes", BVR designated target)
%   * Beam geometry imprecision validated against Dash-34 §2.2.4.4.2.1.2
% - File renamed: section-C4-S3-tms-aa-review-2026-02-16.tex → section-C4-S3-tms-aa-review-2026-02-17.tex
%
% Previous Changes (2026-02-16):
% - C4:S3:S5 Introduction rewritten: improved flow, removed redundancy
% - Parágrafo 1: TMS Left as protagonist, SCAN/LOS introduced, cross-refs to S1/S2
% - Parágrafo 2: BVR/ROE context + IFF protocol + symbology (green/yellow)
% - Parágrafo 3: AIFF availability warning + requirements + cross-refs to S5/Dash-34
% - Removed duplicate content between paragraphs 3 and 4
% - Corrected "block configuration" claim (Dash-34 does not provide list)
% - Changed cross-ref S5 from "detailed availability" to "identification procedures"
% - Fixed grammar: "across all F-16 blocks and configurations"
% - Validated alignment with bullets from C4:S3:S5:S5 (AIFF availability)
% - C4:S3:S5 Introduction marked FINAL, ready for integration
%
% Known Issues:
% - C4:S3:S5 subsubsections S3-S5: Narrative structure exists, content incomplete
% - Subsections C4:S3:S1-S4, S6: Narrative structure exists, content incomplete
% - All HOTAS tables: Placeholders only, require full population
% - TGP A-A (S6): Awaits TGP equipment note integration in C2:S1:S3
%
% Next Steps:
% 1. Populate narrative for C4:S3:S5:S3 (TMS IFF Interrogation Usage Table) — hotastable format
% 2. Populate narrative for C4:S3:S5:S4 (Operational Constraints)
% 3. Populate narrative for C4:S3:S5:S5 (AIFF Availability)
% 4. Repeat for remaining subsections S1-S4, S6 and corresponding tables
% 5. Final validation of all cross-references before integration
%
% Integration Notes:
% - This WIP will be integrated simultaneously with:
%   * section-C2-S1-S3-soi-context-dependency-final-2026-02-15.tex
% - Target version: v0.5.0.0 (TMS Air-to-Air chapter)
% - C4:S3 Introduction + C4:S3:S5 Introduction + C4:S3:S5:S1 + C4:S3:S5:S2 ready; remaining content in development

%============================================================================
% TEMPLATE STRUCTURE — Replace section/subsection headers with your content
%============================================================================

%============================================================================
% CONTENT BEGINS HERE (use \chapter{}, \section{}, \subsection{}, etc.)
%============================================================================
%--------------------------------------------------------------------------
% SECTION 4.3: TMS IN AIR-TO-AIR
%--------------------------------------------------------------------------
\section{TMS in Air-to-Air}
\label{sec:C4-S3}

The Target Management Switch (TMS) in air-to-air operations provides hands-on control for target designation, track management, and sensor cueing during weapons employment. The pilot engaged in beyond-visual-range (BVR) intercepts or within-visual-range (WVR) maneuvering cannot divert attention to MFD interaction or remove hands from the flight stick. TMS commands deliver immediate control of radar tracking modes, target prioritization, and identification functions without interrupting aircraft control or visual scan.

As established in \secref{sec:C2-S1-S3} and detailed in \secref{sec:C3-S1-S1}, A-A master mode restricts Sensor of Interest (SOI) designation to MFD formats only—specifically, the Fire Control Radar (FCR), Horizontal Situation Display (HSD), and Targeting Pod (TGP). Consequently, TMS commands in A-A are restricted to FCR, HSD, and TGP formats; the HUD, although providing essential cueing (\secref{sec:C3-S1-S3}), does not receive TMS inputs in this master mode. The Helmet Mounted Cueing System (HMCS) provides independent off-boresight target cueing capability and AIM-9 seeker slaving, but TMS inputs route through the SOI-designated MFD sensor, not through HMCS or HUD symbology.

TMS behavior in air-to-air varies by sensor mode and tracking state. The primary tactical employment of TMS occurs through FCR tracking modes, which manage the complete engagement cycle from search to weapons delivery. Beyond FCR tracking, TMS Left commands IFF (Identification Friend or Foe) interrogation, critical for Rules of Engagement (ROE) compliance before BVR weapons release, and TMS with TGP as SOI manages visual identification and backup tracking independent of radar state. The FCR provides four primary tracking modes, each optimized for specific tactical contexts:

\begin{itemize}
    \item \textbf{Situation Awareness Mode (SAM)} (\dashref{2.3.1.5.9}) tracks one or two targets while maintaining search coverage, balancing situational awareness with engagement capability.
    
    \item \textbf{Single Target Track (STT)} (\dashref{2.3.1.5.8}) dedicates the radar to a single target with maximum accuracy, but continuously radiates and reveals the F-16's position to threat receivers.
    
    \item \textbf{Track-While-Scan (TWS)} (\dashref{2.3.1.5.7}) automatically builds and maintains track files on up to ten targets while continuing to search, optimizing multi-bogey BVR engagements.
    
    \item \textbf{Air Combat Maneuvering (ACM)} (\dashref{2.3.1.5.10}--\dashref{2.3.1.5.12}) modes provide four automated scan patterns (BORE, 10$\times$60, 30$\times$20, SLEW) optimized for close-range visual acquisition with rapid lock and break-lock cycles.
\end{itemize}

TMS commands manage not only tactical functions within each mode, but also transitions between modes. The pilot transitions between SAM, STT, TWS, and search modes (RWS) using TMS inputs, with the resulting mode depending on the current radar state and targeting context. The same physical TMS input produces different results based on whether a target is bugged, whether the radar is in single-target or multi-target tracking, and whether the pilot is prosecuting a primary threat or re-establishing search coverage. This state-dependency requires the pilot to maintain awareness of current radar mode and tracking status to predict TMS behavior accurately.

The following subsections detail TMS behavior in each air-to-air tracking context. \secref{sec:C4-S3-S1} through \secref{sec:C4-S3-S4} examine FCR modes (SAM, STT, TWS, ACM), \secref{sec:C4-S3-S5} covers IFF interrogation, and \secref{sec:C4-S3-S6} addresses TGP air-to-air tracking. Each subsection includes detailed HOTAS tables documenting specific TMS inputs, effects, Dash-34 cross-references, and recommended training missions for hands-on practice.

\textbf{Note:} TMS with HSD for datalink target management (PDLT designation) is covered in \secref{sec:C4-S2-S1}. While PDLT is an air-to-air function, HSD is primarily a situational awareness display rather than a direct targeting sensor, and its TMS procedures differ from FCR, TGP, and IFF operations documented in this section.


%--------------------------------------------------------------------------
% SUBSECTION 4.3.1: TMS IN SAM (SITUATION AWARENESS MODE)
%--------------------------------------------------------------------------
\subsection{TMS in SAM (Situation Awareness Mode)}
\label{sec:C4-S3-S1}

Situation Awareness Mode (SAM) tracks one or two targets (ST SAM or DT SAM) while maintaining search coverage. TMS commands in SAM manage track transitions and target prioritization.

\paragraph{TMS Up in SAM:}
\begin{itemize}
    \item \textbf{On bugged target:} Transition SAM → STT (Single Target Track). Search coverage ceases; FCR dedicates all resources to tracking the bugged target.
    \item \textbf{On search target (DT SAM):} Designate second target as secondary SAM target, entering Dual Target SAM (DT SAM). Radar tracks both targets while searching.
\end{itemize}

\paragraph{TMS Down in SAM (Return-to-Search):}
\begin{itemize}
    \item \textbf{ST SAM:} Drops bugged target, returns to RWS search mode. Target becomes untracked search symbol.
    \item \textbf{DT SAM:} Drops secondary target, returns to ST SAM with primary target still bugged.
\end{itemize}

\paragraph{TMS Right in SAM (short press <1s):}
\begin{itemize}
    \item \textbf{DT SAM:} Toggles bug between primary and secondary targets. Allows pilot to prioritize which target is primary without re-designation.
\end{itemize}

\paragraph{TMS Right in SAM (held >1s):}
Transitions SAM → TWS. Bugged target becomes TWS primary target. Radar begins building track files on other search targets.

\paragraph{Operational context:}
SAM is ideal for Beyond Visual Range (BVR) engagements where the pilot needs track data on one or two targets while maintaining search awareness. TMS Up/Down provides rapid transitions between search, SAM, and STT as the tactical picture evolves.

See \tabref{tab:C4-S3-SAM} for detailed SAM TMS actions.

%--------------------------------------------------------------------------
% SUBSECTION 4.3.2: TMS IN STT (SINGLE TARGET TRACK)
%--------------------------------------------------------------------------
\subsection{TMS in STT (Single Target Track)}
\label{sec:C4-S3-S2}

Single Target Track (STT) dedicates the FCR to tracking a single airborne target with maximum accuracy. All search processing ceases. TMS commands in STT are limited to breaking lock and returning to previous modes.

\paragraph{TMS Down in STT (Return-to-Search):}
\begin{itemize}
    \item \textbf{First press:} STT → SAM (or DT SAM if secondary target was being extrapolated). STT target becomes SAM primary target.
    \item \textbf{Second press:} SAM → RWS. Bugged target becomes untracked search symbol.
\end{itemize}

\paragraph{TMS Up, TMS Right, TMS Left in STT:}
\begin{itemize}
    \item \textbf{ACM STT:} TMS Up, TMS Right, or CURSOR slew do NOT break lock. This allows pilot to slew scan volume or step targets without losing track (ACM-specific behavior).
    \item \textbf{Non-ACM STT:} TMS Up/Right/Left have no effect. Only TMS Down breaks lock.
\end{itemize}

\paragraph{Operational note:}
STT reveals the F-16's position to any target with RWR (Radar Warning Receiver) due to continuous radar emissions. Use STT only when weapon parameters are achieved and giving away position is acceptable. For covert BVR engagements, SAM or TWS is preferable.

See \tabref{tab:C4-S3-STT} for STT TMS actions.

%--------------------------------------------------------------------------
% SUBSECTION 4.3.3: TMS IN TWS (TRACK-WHILE-SCAN)
%--------------------------------------------------------------------------
\subsection{TMS in TWS (Track-While-Scan)}
\label{sec:C4-S3-S3}

Track-While-Scan (TWS) automatically builds and maintains track files on up to 10 targets while continuing to search. TMS commands manage target prioritization (bugging) and transitions to STT for weapons employment.

\paragraph{TMS Up in TWS:}
\begin{itemize}
    \item \textbf{Cursor over track file:} Bug that target (make it primary/TOI). If already bugged, transition TWS → STT on bugged target.
    \item \textbf{Cursor over search target:} Designate search target, build track file, and bug it.
\end{itemize}

\paragraph{TMS Down in TWS (Return-to-Search):}
\begin{itemize}
    \item \textbf{First press:} Drop bug (if bugged target exists). Radar continues TWS with all track files.
    \item \textbf{Second press:} Dump all track files, return to RWS search mode.
    \item \textbf{Third press:} (from RWS after two presses) Radar rebuilds tracks automatically.
\end{itemize}

\paragraph{TMS Right in TWS (short press <1s):}
Steps bug to next highest-priority track file. Priority is determined by range and track file order.

\paragraph{TMS Right in TWS (held >1s):}
\begin{itemize}
    \item \textbf{With bugged target:} TWS → DT SAM. Bugged target becomes SAM primary; highest-priority track file becomes SAM secondary.
    \item \textbf{Without bugged target:} TWS → RWS (last exited CRM search mode).
\end{itemize}

\paragraph{TMS Up (held >1s) in TWS:}
Initiates spotlight search (same as RWS). Spotlight scan centered on ACQ cursor; track updates on TOI and AIM-120 in-flight targets occur only if spotlight volume includes them.

\paragraph{Operational context:}
TWS excels in multi-target environments (BARCAP, HAVCAP, large force engagements). Pilot can maintain SA on 8-10 threats while bugging and prosecuting the highest-priority target. TMS Right stepping allows rapid target prioritization without MFD interaction.

See \tabref{tab:C4-S3-TWS} for TWS TMS actions.

%--------------------------------------------------------------------------
% SUBSECTION 4.3.4: TMS IN ACM (AIR COMBAT MANEUVERING) MODES
%--------------------------------------------------------------------------
\subsection{TMS in ACM (Air Combat Maneuvering) modes}
\label{sec:C4-S3-S4}

Air Combat Maneuvering (ACM) modes provide automated close-range target 
acquisition (BORE, 10$\times$60, 30$\times$20, SLEW). TMS commands manage 
mode transitions, target rejection, and HMCS integration.

\subsubsection*{TMS Up in ACM}
\label{sec:C4-S3-S4-S1}

\paragraph{ACM BORE with HMCS (TMS Up held):}
\begin{itemize}[nosep]
    \item \textbf{Hold phase (TMS Up depressed):}
    \begin{itemize}[nosep]
        \item FCR scan slaves to HMCS aiming cross position.
        \item Radar in NO RAD state (non-radiating) --- pilot aims scan volume with head movement.
        \item ACQ cursor tracks HMCS aiming cross.
        \item No minimum hold duration --- pilot controls timing.
    \end{itemize}
    
    \item \textbf{Release phase (TMS Up released):}
    \begin{itemize}[nosep]
        \item FCR radiates at HMCS-designated position.
        \item Auto-acquisition attempts on target under aiming cross.
        \item If successful: Enters STT, solid TD box displayed on HUD.
        \item If unsuccessful: Remains in ACM BORE, radiating.
    \end{itemize}
    
    \item \textbf{Tactical use:}
    \begin{itemize}[nosep]
        \item High off-boresight (HOBS) acquisition: Lock target 60--90$^\circ$ off nose without 
              pointing aircraft.
        \item Visual merge: Identify bandit visually, slave radar via HMCS, acquire without 
              nose alignment (critical in scissors/rolling fights).
        \item Target sorting: Acquire specific bandit in multi-bogey merge without broad 
              scan hitting wrong target.
    \end{itemize}
\end{itemize}

\paragraph{ACM BORE without HMCS (TMS Up held):}
\begin{itemize}[nosep]
    \item Holding TMS Up inhibits auto-acquisition until released.
    \item Allows pilot to position boresight cross over desired target in clustered 
          group before commanding acquisition (manual target sorting).
    \item No HMCS slaving --- radar remains boresight-aligned.
\end{itemize}

\paragraph{ACM STT:}
\begin{itemize}[nosep]
    \item TMS Up, TMS Right, or CURSOR slew does NOT break lock (ACM-specific behavior).
    \item Allows pilot to reposition scan or step targets without losing track.
\end{itemize}

\subsubsection*{TMS Down in ACM}
\label{sec:C4-S3-S4-S2}

\paragraph{Return-to-Search (any radiating ACM mode):}
\begin{itemize}[nosep]
    \item First press: Drops radar track (if any), commands 30$\times$20 \texttt{NO RAD}.
    \item Second press: Switches to 10$\times$60 radiating.
    \item Third press: Returns to 30$\times$20 \texttt{NO RAD}.
\end{itemize}

\paragraph{HUD as SOI with FCR slaved to HMCS:}
\begin{itemize}[nosep]
    \item TMS Down commands ACM BORE NO RAD.
\end{itemize}

\subsubsection*{TMS Right in ACM}
\label{sec:C4-S3-S4-S3}

\begin{itemize}[nosep]
    \item Cycles through ACM scan patterns (implementation varies by mode).
    \item Typical sequence: 30$\times$20 $\rightarrow$ SLEW $\rightarrow$ 10$\times$60 $\rightarrow$ BORE.
\end{itemize}

\subsubsection*{Operational context}
\label{sec:C4-S3-S4-S4}

ACM modes automate acquisition in visual merge/dogfight. TMS Up/Down provide 
rapid lock/break-lock cycles essential in turning fights. HMCS integration 
enables HOBS acquisition without pointing aircraft nose.

\textbf{Reference:} ACM modes details in \dashref{2.3.1.5.10}. HMCS 
integration in \dashref{2.5}.

See \tabref{tab:aa-acm-tms} for ACM TMS actions.

%--------------------------------------------------------------------------
% SUBSECTION 4.3.5: TMS AND IFF INTERROGATION
%--------------------------------------------------------------------------

\subsection{TMS and IFF Interrogation}
\label{sec:C4-S3-S5}

Identification Friend or Foe (IFF) interrogation is an air-to-air identification function commanded by \textbf{TMS Left}, enabling the pilot to distinguish friendly aircraft from unknowns and bogeys without removing hands from the flight stick. The interrogation mode executed depends on press duration: a \textit{short press activates SCAN interrogation} for broad area coverage, while a \textit{long press activates Line of Sight (LOS) interrogation} for precise single-target queries. Both modes display interrogation responses on the Fire Control Radar (FCR) MFD with symbology overlaid on radar contacts, enabling immediate threat discrimination during multi-target engagements without requiring MFD interaction or datalink correlation (see \secref{sec:C4-S3-S5-S1} and \secref{sec:C4-S3-S5-S2} for details).

IFF interrogation is operationally critical in beyond-visual-range (BVR) engagements, where visual identification is impossible and Rules of Engagement (ROE) require positive identification before weapons employment to prevent fratricide. The interrogation process begins when the interrogating aircraft transmits a coded query to the target aircraft, whose transponder validates the query and transmits a response back to the interrogator. The FCR then displays identification symbology based on the response received.

Not all F-16 variants are equipped with IFF interrogation capability, as the AIFF (Advanced Interrogator Friend or Foe) system is not universally installed across all F-16 blocks and configurations. When available, TMS Left IFF interrogation requires the AIFF system to be powered and configured via the Up Front Controller (UFC) or backup IFF panel controls. Pilots must verify AIFF availability via cockpit panel configuration before flight; see \secref{sec:C4-S3-S5-S5} for cockpit identification methods and block/variant differences. For complete IFF system architecture, operational modes, and control panel operation, see \dashref{2.2.4}.

%--------------------------------------------------------------------------
% SUBSUBSECTION 4.3.5.1: SCAN INTERROGATION MODE
%--------------------------------------------------------------------------
\subsubsection{SCAN Interrogation Mode (TMS Left Short Press)}
\label{sec:C4-S3-S5-S1}

SCAN interrogation mode interrogates multiple radar contacts simultaneously to identify friendly and hostile aircraft. A \textbf{short press of TMS Left (less than 0.6 second)} activates SCAN mode, commanding the AIFF subsystem to transmit interrogation queries and display identification symbology on the FCR. The system can display up to 32 target replies; if more than 32 aircraft respond, only the closest 32 are reported based on range.

\paragraph{Configuration:} SCAN coverage can be configured via \texttt{DED INTG} page (\texttt{LIST RCL} button) or FCR CNTL page OSB 10 to decouple from or couple to the radar scan volume; see \dashref{2.2.4.4.2.1.1} and \dashref{2.3.1.5.1.6.2}.

The FCR displays target replies as symbology overlaid on radar contacts, enabling immediate visual discrimination without requiring MFD interaction. A green circle identifies a friendly aircraft that provided a correct response to the interrogation, while a yellow square identifies a bogey or unknown contact that failed to respond or provided an incomplete response.

TMS Left short press is employed when the pilot requires identification of multiple contacts across the tactical picture. This is common in beyond-visual-range (BVR) intercepts where the radar displays several unknown tracks and Rules of Engagement (ROE) require positive identification before weapons employment.

%--------------------------------------------------------------------------
% SUBSUBSECTION 4.3.5.2: LOS INTERROGATION MODE
%--------------------------------------------------------------------------
\subsubsection{LOS Interrogation Mode (TMS Left Long Press)}
\label{sec:C4-S3-S5-S2}

Line of Sight (LOS) interrogation mode interrogates targets within the \textit{beam coverage area centered around the acquisition cursor}, providing precise cursor-directed interrogation when the pilot requires identification of contacts at a specific cursor position rather than across the radar scan volume. Unlike SCAN mode, which interrogates multiple targets across the radar scan volume (see \secref{sec:C4-S3-S5-S1}), LOS concentrates interrogation resources on the beam area defined by the cursor placement, using two beams centered on the cursor to provide coverage while minimizing interrogation of contacts outside the cursor-designated area.

A \textbf{long press of TMS Left (at least 0.6 second)} activates LOS mode, commanding the AIFF subsystem to interrogate the target at the acquisition cursor position. The system can display up to 16 target replies, compared to 32 in SCAN mode, reflecting the narrower interrogation focus. The interrogation beams are positioned to provide optimal coverage around the cursor, though due to fixed beam geometry the beams may not align precisely with the exact cursor position. Responses are filtered by range, starting from the cursor's current range position and extending to the maximum selected range on the MFD, ensuring that only relevant targets within the designated range bracket are displayed.

\paragraph{Configuration:} LOS mode is configured via the \texttt{DED INTG} page (\texttt{LIST RCL} button, then \texttt{DCS} right to sequence to the LOS page), which operates independently from the SCAN page and allows the pilot to set different active modes and codes for each interrogation type. See \dashref{2.2.4.4.2.1} for complete \texttt{INTG} page operation.

The FCR displays target replies using the same symbology as SCAN mode (see \secref{sec:C4-S3-S5-S1}): green circle for friendly, yellow square for bogey/unknown. See \dashref{2.2.4.4.2.1.2} for complete LOS interrogation behavior and \dashref{2.2.4.4.2.2} for response symbology details.

TMS Left long press (LOS) is employed when the pilot requires identification at a specific cursor position rather than across the full scan volume. This is common when prosecuting a designated target in beyond-visual-range engagements.

%--------------------------------------------------------------------------
% SUBSUBSECTION 4.3.5.3: TMS IFF INTERROGATION USAGE TABLE
%--------------------------------------------------------------------------
\subsubsection{TMS IFF Interrogation Usage Table}
\label{sec:C4-S3-S5-S3}

% PARÁGRAFO — Table introduction (2-3 sentenças)
% - The table below summarizes TMS Left behavior in IFF interrogation contexts across SCAN and LOS modes
% - Each entry specifies press duration, interrogation mode activated, system behavior, Dash-34 references, and training missions

% HOTASTABLE — tab:C4-S3-IFF
% Colunas: Mode | Dir. | Act. | Sensor/Mode | Effect/Nuance | Dash-34 | Training
% LINHA 1:
% Mode: A-A
% Dir.: Left
% Act.: Short (<0.6s)
% Sensor/Mode: AIFF SCAN
% Effect/Nuance: Activates SCAN interrogation mode. In DCPL: ±60° full scan (8 beam positions, up to 32 replies). In CPL: scan volume coupled to FCR field-of-regard. Green circle = friendly (correct response), yellow square = bogey/unknown (incomplete/no response). Replies displayed 2 seconds per mode. CPL/DCPL configured via DED SCAN INTG page (LIST RCL) or FCR CNTL page OSB 10.
% Dash-34: 2.2.4.4.2.1.1, 2.2.4.4.2.2, 2.3.1.5.1.6.2
% Training: 18

% LINHA 2:
% Mode: A-A
% Dir.: Left
% Act.: Long (≥0.6s)
% Sensor/Mode: AIFF LOS
% Effect/Nuance: Activates LOS interrogation mode. Interrogates single target along FCR LOS (bugged target in SAM/STT, cursor position in RWS). Up to 16 replies displayed (closest 16 if >16 respond). Range filtered from cursor to max MFD range; no azimuth/elevation filtering. Precise single-target interrogation. Replies displayed 1.5 seconds per mode. Cursor position not updated during multi-mode sequence.
% Dash-34: 2.2.4.4.2.1.2, 2.2.4.4.2.2
% Training: 18

%--------------------------------------------------------------------------
% SUBSUBSECTION 4.3.5.4: OPERATIONAL CONSTRAINTS AND NOTES
%--------------------------------------------------------------------------
\subsubsection{Operational Constraints and Notes}
\label{sec:C4-S3-S5-S4}

% PARÁGRAFO 1 — TMS/cursor interaction constraints (3-4 sentenças)
% - During LOS multi-mode interrogation sequence, cursor position not updated between modes being interrogated
% - If cursor moved significantly during interrogation, responses may not appear near currently displayed cursor position
% - AIFF continues interrogating at initial cursor position captured at start of sequence
% - Pilots should stabilize cursor on target before initiating LOS interrogation for accurate response correlation
% - In SCAN mode, this constraint does not apply (SCAN interrogates predefined area, not cursor-referenced position)

% PARÁGRAFO 2 — Radar mode flexibility (3-4 sentenças)
% - LOS interrogation can be initiated in any A-A radar mode, including standby or off states
% - SCAN in CPL mode requires active radar mode to couple AIFF FOV to FCR field-of-regard
% - SCAN in DCPL mode operates independently of radar state, can be used with radar in standby
% - CPL/DCPL configuration persists across radar mode changes until manually altered via DED or FCR CNTL page
% - Pilots should verify DCPL/CPL configuration before flight (DCPL for broad area coverage, CPL for radar-focused interrogation)

%--------------------------------------------------------------------------
% SUBSUBSECTION 4.3.5.5: AIFF AVAILABILITY AND BLOCK DIFFERENCES
%--------------------------------------------------------------------------
\subsubsection{AIFF Availability and Block Differences}
\label{sec:C4-S3-S5-S5}

% PARÁGRAFO [X] — AIFF availability and variant differences (8-10 sentenças)
% - TMS Left IFF interrogation requires aircraft equipped with AN/APX-113(V) AIFF system
% - AIFF integrated into modern F-16 variants as internal avionics (not external pod)
% - AIFF-equipped aircraft feature IFF control panel in cockpit with backup IFF controls
% - Non-AIFF aircraft retain older AUX COMM panel with backup TACAN controls
% - Not all F-16 variants possess interrogator capability
% - All aircraft have IFF transponders (respond to interrogations from other aircraft/ground stations)
% - Only AIFF-equipped aircraft can actively interrogate other aircraft via TMS Left
% - In non-AIFF aircraft, TMS Left has no function — pilot cannot initiate hands-on IFF interrogation
% - The Dash-34 does not provide a definitive list of AIFF-equipped blocks
% - Pilots must verify cockpit panel configuration before flight to determine interrogator availability:
%   → IFF panel = AIFF capability (interrogator + transponder)
%   → AUX COMM panel = transponder-only configuration (no interrogator)
% - See \dashref{2.2.4.2} for AIFF system architecture, beam forming network, antenna configuration
% - See \dashref{2.2.4.3} for IFF control panel operation, IFF MASTER switch, backup controls

PARÁGRAFO [X]: AIFF vs Link 16 context
- Aircraft without AIFF rely on Link 16 or AWACS for target ID
- AIFF enables hands-on interrogation independent of datalink
- Critical when: (a) Link 16 unavailable, (b) tracks not correlated via datalink
- Dash-34: [identificar seção relevante se houver]
- Rationale: Explains why AIFF matters even with modern datalink capabilities

%--------------------------------------------------------------------------
% SUBSECTION 4.3.6: TMS with Targeting Pod in Air-to-Air Mode
%--------------------------------------------------------------------------
\subsection{TMS with Targeting Pod in Air-to-Air Mode}
\label{sec:C4-S3-S6}

The Targeting Pod (TGP) operates in Air-to-Air mode for visual/IR target 
identification and tracking independent of FCR. TMS commands with TGP as 
SOI manage POINT track designation and track rejection.

\subsubsection*{TMS Up (TGP SOI, A-A mode)}
\label{sec:C4-S3-S6-S1}

\begin{itemize}[nosep]
    \item Commands POINT track on center of TGP field-of-view (crosshair position).
    \item TGP locks onto visual/IR target at designated position.
    \item After track command, TGP LOS becomes independent of FCR LOS.
    \item HUD displays:
    \begin{itemize}[nosep]
        \item Dotted 50-mr TD box (TGP track position).
        \item Solid 50-mr TD box (FCR track position) if FCR also tracking.
        \item Dotted TLL (Target Locator Line) with ``T'' + angle if TGP track outside HUD FOV.
    \end{itemize}
\end{itemize}

\subsubsection*{TMS Down (TGP SOI, A-A mode)}
\label{sec:C4-S3-S6-S2}

\begin{itemize}[nosep]
    \item Breaks TGP POINT track.
    \item Returns TGP to slave mode --- TGP re-slaves to FCR LOS.
    \item Dotted TD box/TLL removed from HUD.
\end{itemize}

\subsubsection*{Operational context}
\label{sec:C4-S3-S6-S3}

TGP A-A mode provides:
\begin{enumerate}[nosep]
    \item \textbf{Visual ID (VID):} Identify bandits before BVR shot (ROE compliance).
    \item \textbf{Dual tracking:} FCR tracks primary threat, TGP tracks secondary (2v1 scenarios).
    \item \textbf{Backup tracking:} If FCR loses lock (chaff/jamming), TGP maintains visual track.
    \item \textbf{High-aspect targets:} TGP may track targets FCR cannot resolve (low RCS, 
          beam aspect).
\end{enumerate}

\subsubsection*{Limitations}
\label{sec:C4-S3-S6-S4}

\begin{itemize}[nosep]
    \item TGP A-A track is passive (no ranging) --- uses triangulation for range estimate.
    \item Track quality degrades at long range ($>$10~NM) or high closure rates.
    \item TGP slew rate slower than FCR --- not suitable for high-G dogfight tracking.
\end{itemize}

\subsubsection*{Procedure (VID example)}
\label{sec:C4-S3-S6-S5}

\begin{enumerate}[nosep]
    \item FCR in RWS/SAM with unknown contact.
    \item Select A-A master mode, TGP powered.
    \item DMS Down to make TGP SOI.
    \item Slew TGP crosshair over bandit (TGP initially slaved to FCR LOS).
    \item TMS Up --- TGP commands POINT track, dotted TD box appears on HUD.
    \item Visually identify via TGP video (MFD TGP page).
    \item TMS Down --- break TGP track after ID complete.
\end{enumerate}

\textbf{Reference:} TGP A-A mode in \dashref{2.6.1.7.5}. SRM procedures 
with TGP in \dashref{4.3.2.4.3.12}.

See \tabref{tab:aa-tgp-tms} for TGP A-A TMS actions.

%--------------------------------------------------------------------------
% PLACEHOLDER: TABLES FOR SECTION 4.3
%--------------------------------------------------------------------------

[PLACEHOLDER: The following tables will be developed with full hotastable format:]

% \subsection{TMS actions in SAM (Table)}
% \begin{table}[H]
% \begin{hotastable}{TMS Actions in SAM (Situation Awareness Mode)}
% [Content: ST SAM, DT SAM | TMS Up/Down/Right | Designate, RTS, bug toggle | Dash-34 §2.3.1.5.9 | TRN 18]
% \end{hotastable}
% \label{tab:C4-S3-SAM}
% \end{table}

% \subsection{TMS actions in STT (Table)}
% \begin{table}[H]
% \begin{hotastable}{TMS Actions in STT (Single Target Track)}
% [Content: STT, ACM STT | TMS Down RTS | Dash-34 §2.3.1.5.8 | TRN 18]
% \end{hotastable}
% \label{tab:C4-S3-STT}
% \end{table}

% \subsection{TMS actions in TWS (Table)}
% \begin{table}[H]
% \begin{hotastable}{TMS Actions in TWS (Track-While-Scan)}
% [Content: TWS | TMS Up/Down/Right | Bug, RTS, step, spotlight | Dash-34 §2.3.1.5.7 | TRN 18]
% \end{hotastable}
% \label{tab:C4-S3-TWS}
% \end{table}

% \subsection{TMS actions in ACM modes (Table)}
% \begin{table}[H]
% \begin{hotastable}{TMS Actions in ACM Modes}
% [Content: ACM BORE, 10x60, 30x20, SLEW | TMS Up/Down/Right | HMCS slave, RTS, mode cycle | Dash-34 §2.3.1.5.10-12 | TRN 20]
% \end{hotastable}
% \label{tab:C4-S3-ACM}
% \end{table}

% \subsection{TMS actions for IFF interrogation (Table)}
% \begin{table}[H]
% \begin{hotastable}{TMS Left — IFF Interrogation}
% [Content: A-A | TMS Left short/long | SCAN/LOS interrogation | Dash-34 IFF chapter | TRN 18]
% \end{hotastable}
% \label{tab:C4-S3-IFF}
% \end{table}

\end{document}